\section{Span, Linear Independence, Basis, and Dimension}

\subsection{Spanning Set and Linear Independence}

\begin{definition}[Span of a Set, Spanning Set]
    Let \(S \subseteq V\). The \term{span} of \(S\) in \(V\) is given by
    \[
        \spn \pqty{S} = \angBrac{S} = \set{{\sum_{\vb{s} \in S}}' \lambda_{\vb{s}} \vb{s}}{\lambda_{\vb{s}} \in \FF},
    \]
    where the sum \({\sum}'\) denotes a \term{finite sum}, i.e., we require that \(\lambda_{\vb{s}} = 0\) for all but finitely many \(\vb{s} \in S\).
    
    The elements of \(\spn\pqty{S}\), in the finite sum form, are called \term{linear combinations} of elements of \(S\).

    If \(\spn\pqty{S} = V\), then we say \(S\) is a \term{spanning set} of \(V\).
\end{definition}

\begin{remark}
    Note that we require \({\sum}'\), since there is no well-defined way to sum infinitely many vectors in a vector space (this requires \emph{analysis}). In particular, if \(S\) is infinite, then \(\spn\pqty{S}\) is the set of all finite linear combinations of elements of \(S\).
\end{remark}

\begin{definition}[Finite/Infinite-Dimensional Vector Space]
    We say that a vector space \(V\) is \term{finite-dimensional} if there exists a finite spanning set \(S \subseteq V\). Otherwise, we say that \(V\) is \term{infinite-dimensional}.
\end{definition}

\begin{remarks}
    \begin{itemize}
        \item \(\spn\pqty{S} \subspace V\) for any \(S \subseteq V\).
        \item By convention, \(\spn\pqty{\emptyset} = \Bqty{\vb{0}_V}\).
        \item If \(S\) spans \(V\) and \(S \subseteq \spn\pqty{T}\), then \(T\) spans \(V\).
        \item \(\spn\pqty{S \cup T} = \spn\pqty{S} + \spn\pqty{T}\).
    \end{itemize}
\end{remarks}

\begin{examples}
    \begin{enumerate}[label=(\roman*)]
        \item Let \(V = \RR^3\), and let
        \[
            S = \Bqty{\pmqty{1 \\ 0 \\ 0}, \pmqty{1 \\ 1 \\ 2}} \qand T = \Bqty{\pmqty{2 \\ 1 \\ 2}, \pmqty{0 \\ 1 \\ 2}, \pmqty{-1 \\ 2 \\ 4}}.
        \]

        Then
        \[
            \spn\pqty{S} = \spn\pqty{T} = \set{\pmqty{x \\ y \\ 2y}}{x, y \in \RR} \subspace \RR^3.
        \]

        \item Let \(V = \RR^3\), and let
        \[
            T = \set{\kronecker_n}{n \in \NN} \subseteq V
        \]
        where \(\functiondomain{\kronecker_n}{\NN}{\RR}\) is given by
        \[
            \kronecker_n\pqty{m} = \kronecker_{n, m} = \begin{cases}
                1 & n = m, \\
                0 & n \neq m.
            \end{cases}
        \]

        Then \(\spn\pqty{T} \neq \RR^\NN\). For instance, if \(f \in \RR^\NN\) satisfies \(f\pqty{n} = 1\) for all \(n\), then \(f \notin \spn\pqty{T}\). In fact,
        \[
            \spn\pqty{T} = \set{g \in V}{g\pqty{n} = 0 \text{ for sufficiently large }n \in \NN}.
        \]
    \end{enumerate}
\end{examples}

\begin{definition}[Linearly Independence/Dependence]
    A set \(S \subseteq V\) is \term{linearly independent (L.I.)} if, for any linear combination satisfying
    \[
        {\sum_{\vb{s} \in S}}' \lambda_{\vb{s}} \vb{s} = \vb{0}_V,
    \]
    we have \(\lambda_{\vb{s}} = 0\) for all \(\vb{s} \in S\). Otherwise, we say that \(S\) is \term{linearly dependent (L.D.)}.
\end{definition}

\begin{remarks}
    \begin{itemize}
        \item If a set \(S \subseteq V\) is linearly independent, then \(\vb{0} \notin S\).
        \item If \(S\) is finite, with \(S = \Bqty{\vb{v}_1, \vb{v}_2, \ldots, \vb{v}_n}\) and \(\abs{S} = n\), then \(S\) is linearly independent, if and only if for all \(\lambda_1, \lambda_2, \ldots, \lambda_n \in \FF\), that
        \[
            \sum_{i = 1}^{n} \lambda_i \vb{v}_i = \vb{0} \implies \lambda_i = \cdots = \lambda_n = 0.
        \]
        \item If \(S\) is infinite, then \(S\) is linearly independent, if and only if every finite subset of \(S\) is linearly independent.
        \item Every subset of a linearly independent set is itself linearly independent.
    \end{itemize}
\end{remarks}

\begin{examples}
    \begin{enumerate}[label=(\roman*)]
        \item Let \(V = \RR^3\), and let
        \[
            S = \Bqty{\pmqty{1 \\ 0 \\ 0}, \pmqty{1 \\ 1 \\ 2}} \qand T = \Bqty{\pmqty{2 \\ 1 \\ 2}, \pmqty{0 \\ 1 \\ 2}, \pmqty{-1 \\ 2 \\ 4}}.
        \]

        Then \(S\) is linearly independent, while \(T\) is linearly dependent, since
        \[
            1 \cdot \pmqty{2 \\ 1 \\ 2} - 5 \cdot \pmqty{0 \\ 1 \\ 2} + 2 \cdot \pmqty{-1 \\ 2 \\ 4} = \pmqty{0 \\ 0 \\ 0}.
        \]

        On the other hand, if \(\lambda, \mu \in \RR\) is such that
        \[
            \vb{0} = \lambda \pmqty{1 \\ 0 \\ 0} + \mu \pqty{1 \\ 1 \\ 2} = \pmqty{\lambda + \mu \\ \mu \\ 2\mu}
        \]
        then \(\lambda = \mu = 0\), so \(S\) is linearly independent.

        \item Let \(V = \RR^3\), and let
        \[
            T = \set{\kronecker_n}{n \in \NN} \subseteq V
        \]
        where \(\functiondomain{\kronecker_n}{\NN}{\RR}\) is given by
        \[
            \kronecker_n\pqty{m} = \kronecker_{n, m} = \begin{cases}
                1 & n = m, \\
                0 & n \neq m.
            \end{cases}
        \]

        Let \(f = {\sum_{n \in \NN}}'\lambda_n \delta_n\) where \(\lambda_n \in \RR\). If \(f = 0\), then for all \(m \in \NN\),
        \[
            0 = f\pqty{m} = \pqty{{\sum_{n \in \NN}' \lambda_n \delta_n}} \pqty{m} = {\sum_{n \in \NN}}' \lambda_n \delta_n\pqty{m} = \lambda_m,
        \]
        precisely giving the trivial linear combination. so \(T\) is linearly independent.
    \end{enumerate}
\end{examples}

\subsection{Basis}

\begin{definition}[Basis]
    Let \(S \subseteq V\). \(S\) is a \term{basis} of \(V\) if it is linearly independent and spans \(V\).
\end{definition}

\begin{examples}
    \begin{enumerate}
        \item Let \(\vb{e}_i \in \FF^n\) be the \term{\(i\)-th standard basis vector}, with entries 0 except for the \(i\)-th entry, which is 1. Then the set \(\Bqty{\vb{e}_i}{1 \leq i \leq n}\) is the \term{standard basis} of \(\FF^n\).
        \item The set
        \[
            \Bqty{\pmqty{1 \\ 0 \\ 0}, \pmqty{1 \\ 1 \\ 0}, \pmqty{1 \\ 1 \\ 1}} \subseteq \FF^3
        \]
        is a basis for \(\FF^3\).
        \item Let \(p_n \in P\pqty{\RR}\) be given by \(p_n\pqty{x} = x^n\). Then the set \(\Bqty{p_n}{n \in \NN}\) is a basis for \(P\pqty{\RR}\).
    \end{enumerate}
\end{examples}

\begin{proposition}[Finite Spanning Set reduces to a Basis]
    Suppose \(S \subseteq V\) is a finite spanning set. Then there exists a subset \(S' \subseteq S\) such that \(S\) is a basis of \(V\).
\end{proposition}

\begin{proof}
    If \(S\) is linearly independent, then we are done. If not, write
    \[
        S = \Bqty{\vb{v}_1, \vb{v}_2, \ldots, \vb{v}_n}.
    \]

    There exists \(\lambda_1, \lambda_2, \ldots, \lambda_n \in \FF\), not all zero, such that
    \[
        \sum_{i = 1}^{n} \lambda_i \vb{v}_i = \vb{0}.
    \]

    Without loss of generality, up to relabelling, that \(\lambda_1 \neq 0\). Then rearranging this expression,
    \[
        \vb{v}_1 = - \frac{1}{\lambda_1} \sum_{i = 2}^{n} \lambda_i \vb{v}_i \in \spn\pqty{S \setminus \Bqty{\vb{v}_1}}.
    \]

    So \(S \setminus \Bqty{\vb{v}_1}\) also spans \(V\). Repeat this and this will terminate after at most \(n\) steps.
\end{proof}

\begin{corollary}[Finite-Dimensional Vector Space has a Finite Basis]
    \label{cor:finite-dimensional-vector-space-has-finite-basis}%
    Every finite-dimensional vector space has a finite basis.
\end{corollary}

\begin{remark}
    We will \emph{not} prove that every vector space has a basis. Some non-finite-dimensional vector spaces do have a basis, for example, the above set of polynomial functions over \(\RR\). However, even attempting to write down a basis of \({\RR}^{\NN}\) will be a lot of trouble -- in fact you cannot write down a basis of this space explicitly. The existence of a basis for every vector space is equivalent to the \emph{axiom of choice}, which is a controversial axiom in set theory. In other words, one should not assume that a infinite-dimensional vector space has a basis, unless it is explicitly constructed.
\end{remark}

\begin{theorem}[Steinitz Exchange Lemma]
    \label{thm:steinitz-exchange-lemma}%
    Let \(V\) be a \(\FF\)-vector space, \(S, T \subseteq V\) finite, with \(S\) being linearly-independent, and \(T\) spanning \(V\). Then we have the following.

    \begin{itemize}
        \item \emph{Linear independence set is always smaller than spanning set.} \(\abs{S} \leq \abs{T}\).
        \item \emph{Removing elements from a spanning set and replacing with linearly independent elements.} There exists \(T' \subseteq T\), such that \(S \cup T'\) spans \(V\), and \(\abs{T'} = \abs{T} - \abs{S}\).
    \end{itemize}
\end{theorem}

\begin{proof}
    Let \(S = \Bqty{\vb{s}_1, \ldots, \vb{s}_m}\) nad \(T = \Bqty{\vb{t}_1, \ldots, \vb{t}_n}\), with \(\abs{S} = m\) and \(\abs{T} = n\).

    If \(S = \emptyset\), then we are done. If not, we conduct the following steps.

    \textbf{Step 1.} Since \(\spn\pqty{T} = V\), then there exists \(\lambda_i \in \FF\), such that
    \[
        \vb{s}_1 = \sum_{i = 1}^{N} \lambda_i \vb{t}_i,
    \]
    and since \(\vb{s}_1 \neq \vb{0}\), so \(\lambda_i\) cannot be trivially all zeros. Without loss of generality, up to relabelling, we can assume that \(\lambda_1 \neq 0\). Then rearranging this expression,
    \[
        \vb{t}_1 = \frac{1}{\lambda_1} \pqty{\vb{s}_1 - \sum_{i = 2}^{N} \lambda_i \vb{t}_i} \in \spn\pqty{\Bqty{\vb{s}} \cup \Bqty{\vb{t}_2, \ldots, \vb{t}_n}}.
    \]

    \textbf{Step \(\boldsymbol{\pqty{k + 1}}\)} where \(k < m\). Assume the set \(\Bqty{\vb{s}_1, \ldots, \vb{s}_k, \vb{t}_{k + 1}}, \ldots, \vb{t}_n\) spans \(V\). Write
    \[
        \vb{s}_{k + 1} = \sum_{i = k + 1}^{n} \lambda_i \vb{t}_i + \sum_{j = 1}^{k} \mu_j \vb{s}_j,
    \]
    where \(\lambda_i, \mu_j \in \FF\), not all \(0\).

    If all \(\lambda_i = 0\), then
    \[
        \vb{s}_{k + 1} = \sum_{j = 1}^{k} \mu_j \vb{s}_j = \vb{0}
    \]
    but \(S\) is linearly independence, giving a contradiction.

    Hence, without loss of generality, \(\lambda_{k + 1} \neq 0\), so we can write
    \[
        \vb{t}_{k + 1} = \frac{1}{\lambda_{k + 1}} \pqty{\vb{s}_{k + 1} - \sum_{i = k + 2}^{n} \lambda_i \vb{t}_i - \sum_{j = 1}^{k} \mu_j \vb{s}_j} \in \spn\pqty{\Bqty{\vb{s}_1, \ldots, \vb{s}_{k + 1}} \cup \Bqty{\vb{t}_{k + 2}, \ldots, \vb{t}_n}}.
    \]

    So the new set \(\Bqty{\vb{s}_1, \ldots, \vb{s}_{k + 1}, \vb{t}_{k + 2}, \ldots, \vb{t}_n}\) spans \(V\).

    We repeat the process until either we run out of elements of \(S\), or elements of \(T\). But if we exhausted \(T\) before \(S\), then we would have a linear combination of elements of \(S\) that gives \(\vb{0}\), contradicting the linear independence of \(S\). Hence, we must have exhausted \(S\) first, and we have \(\abs{S} \leq \abs{T}\), which gives \(m \leq n\), and
    \[
        \Bqty{\vb{s}_1, \ldots, \vb{s}_m, \vb{t}_{m + 1}, \ldots, \vb{t}_n} \subseteq V
    \]
    with \(T' = \Bqty{\vb{t}_{m + 1}, \ldots, \vb{t}_n}\).
\end{proof}

\subsection{Dimension}

\begin{corollary}[Basis of Finite-Dimensional Vector Space]
    Let \(V\) be a finite-dimensional vector space. Then every basis for \(V\) is finite, and any two bases \(B_1, B_2\) of \(V\) satisfy \(\abs{B_1} = \abs{B_2}\).
\end{corollary}

\begin{proof}
    By \autoref{cor:finite-dimensional-vector-space-has-finite-basis}, \(V\) has some finite basis \(B\). Suppose \(B'\) is an infinite basis.

    Let \(B'' \subset B'\) with \(\abs{B''} = \abs{B} + 1\). Then \(B''\) is linearly-independent, and \(B\) spans \(V\), contradicting with the first part of \autoref{thm:steinitz-exchange-lemma}.

    If \(B_1, B_2\) are bases of \(V\), then both must be finite. By the first part of \autoref{thm:steinitz-exchange-lemma}, we have \(\abs{B_1} \leq \abs{B_2}\) and \(\abs{B_2} \leq \abs{B_1}\), so \(\abs{B_1} = \abs{B_2}\).
\end{proof}

\begin{definition}[Dimension]
    For a finite-dimensional vector space \(V\), its \term{dimension} \(\dim\pqty{V}\) is the common size of any basis of \(V\).
\end{definition}